% Checked deduction: author, root, vandiver and archive whole-proof readings completed.
% Candidate and first independent audit: work/irregular37_pole_candidate.md.
\subsection{A nonlinear separation argument at \texorpdfstring{$37$}{37}}
\label{sec:kernel37-poles}

The failure of integral odd generation at $37$ does not decide the kernel-field question. The following argument detects its extra even arithmetic direction through a pole in a quadratic moment function. It uses the published cup-product computation, but supplies a separate group-theoretic and analytic proof of faithfulness. The proof is a deduction developed in this run, not a theorem located in the literature or a consequence of step~17.

\begin{theorem}[K17-37-KERNEL]\label{thm:kernel37-poles}
For the maximal pro-$37$ extension $\Omega_{37}$ of $\mathbb Q(\mu_{37})$ unramified outside $37$, the tripod kernel field satisfies
\[
 \Omega_{37}^*=\Omega_{37}.
\]
This is compatible with $(\mathfrak g_{37}^{\mathrm{ab}})_{32}\ne0$ and failure of integral odd generation.
\end{theorem}

We first prove the separation statement needed for the arithmetic application.

\paragraph{K17-POLE-ONE: the compact orbit basis and separation criterion.}
Let $p$ be odd, let $\Gamma=\langle\gamma\rangle\simeq\mathbb Z_p$, and put $\Lambda=\mathbb Z_p[[\Gamma]]$. Suppose $H$ is free pro-$p$ with a continuous $\Gamma$-action, and that its abelianization has the compatible decomposition
\[
 H^{\mathrm{ab}}=
 \left(\bigoplus_{a\in A}\Lambda e_a\right)\oplus\mathbb Z_p e_y,
 \qquad \gamma e_y=\alpha e_y,\qquad \alpha\in1+p\mathbb Z_p.       \tag{N.1}
\]
Here $A$ consists of all odd residue classes modulo $p-1$. A commuting finite group $\Delta$ of order $p-1$ acts on $e_a$ by character $a$ and on $e_y$ by an even character. Choose lifts $x_a,y\in H$, and set
\[
 x_{a,t}=\gamma^t x_a\gamma^{-t},\qquad t\in\mathbb Z_p.
\]
These are a free compact generating space for $H$, in the following precise sense.

Let $F_X$ be the inverse limit of the finite-rank free pro-$p$ groups on
\[
 X_n=(A\times\mathbb Z/p^n)\amalg\{y\},
\]
with maps folding the residue indices. This limit has the universal property for continuous maps from $X=(A\times\mathbb Z_p)\amalg\{y\}$ into a pro-$p$ group: after composition with any finite quotient of the target, such a map factors through an $X_n$. Restricting that property to abelian targets gives
\[
 F_X^{\mathrm{ab}}=\mathbb Z_p[[X]]=
 \Lambda^{|A|}\oplus\mathbb Z_p.
\]
Thus the natural map $F_X\to H$ is an isomorphism on abelianizations. It is onto: a proper closed subgroup of a pro-$p$ group lies in a proper open subgroup and then in a maximal open subgroup, so cannot surject on the Frattini quotient. Since $H$ is free pro-$p$, projectivity supplies a section. Its closed image surjects on $F_X/\Phi(F_X)$ and must therefore be all of $F_X$. The map is an isomorphism. This argument uses neither infinite-rank Hopficity nor an assumed interchange of abelianization and inverse limits.

The compatible finite-rank Magnus embeddings consequently give an injective measure-valued Magnus expansion of $H$. Assign degree $1$ to the orbit letters and degree $2$ to $y$. Every nontrivial element has a first finite nonzero degree. In any degree, each word pattern has a bounded $\mathbb Z_p$-measure as coefficient in the orbit parameters of its letters of degree $1$.

\paragraph{A quadratic meromorphic moment.}
Suppose there is a continuous dense map $\rho:H\to U(\mathbb Q_p)$ into the actual rational tripod image, whose Lie algebra is the completed free Lie algebra on odd generators $s_m$, $m\ge3$. Suppose $\gamma$ is represented by a Levi element with cyclotomic value $u$, a topological generator of $1+p\mathbb Z_p$. Thus it scales $s_m$ by $u^m$, and $\Delta$ acts by character $m$.

Let $\ell_m$ be the $m$th abelian coordinate. Equivariance gives
\[
 \ell_m(\rho(x_{a,t}))=
 \begin{cases}a_m u^{mt},&m\equiv a\pmod{p-1},\\0,&\text{otherwise}.
 \end{cases}
\]
Each $a_m$ is nonzero by density and the compact generating space: $y$ has no odd linear coordinate by parity. Rescale each free letter by $a_m$, weight by weight over $\mathbb Q_p$. The linear part of $\rho(x_{a,t})-1$ now has coefficient $u^{mt}$ on its matching letters. No uniformly integral change of lattice is asserted.

Set
\[
 E=\gamma y\gamma^{-1}y^{-\alpha}\in\overline{[H,H]}.
\]
Its degree-two coefficients involving two orbit letters give bounded measures $\mu_{a,b}$ and transforms
\[
 C_{a,b}(T_1,T_2)=
 \int_{\mathbb Z_p^2}(1+T_1)^t(1+T_2)^v\,d\mu_{a,b}(t,v)
 \in\mathbb Z_p[[T_1,T_2]].
\]
Assume that for one ordered pair $(a_0,b_0)$,
\[
 C_{a_0,b_0}(0,0)\in\mathbb Z_p^\times.                         \tag{N.2}
\]
We will verify this hypothesis arithmetically at $37$.

Let $B_{m,n}$ be the coefficient of $s_ms_n$ in $\rho(y)-1$. Since the latter has no length-one part, comparison in length two yields
\[
 (u^{m+n}-\alpha)B_{m,n}
   =C_{a,b}(u^m-1,u^n-1),\qquad a\equiv m,\ b\equiv n\pmod{p-1}.
                                                                    \tag{N.3}
\]
Indeed the two sides of $E=\gamma y\gamma^{-1}y^{-\alpha}$ give the factor $u^{m+n}-\alpha$, while only the two-orbit-letter Magnus terms can contribute to target length two. Terms of $\overline{[H,H]}$ involving $y$ have assigned degree at least three.

Write
\[
 D(T_1,T_2)=(1+T_1)(1+T_2)-\alpha.
\]
Equation (N.3) represents the length-two kernel of $y$ by $C_{a,b}/D$. Divisions will only be performed in a fraction field after clearing denominators. The numerator in (N.2) is a unit. The divisor is genuine in the open unit polydisc, since $T_2=\alpha/(1+T_1)-1$ and $\alpha-1\in p\mathbb Z_p$.

\paragraph{Ordered tilings and pole separation.}
Suppose $g\ne1$ belongs to $\ker\rho$, and let $d$ be its first nonzero assigned Magnus degree. A source word of degree $d$ specifies a subset
\[
 J\subseteq\{1,\ldots,d-1\},\qquad
 \text{no two elements of }J\text{ adjacent}.
\]
Each $j\in J$ denotes one $y$ filling target positions $j,j+1$; the other positions are singleton orbit letters. For each pattern and its singleton types, let $A_J$ be the Amice transform of its coefficient measure. It depends only on the unpaired variables.

Fix the residue classes of an ordered target word of length $d$. Its vanishing coefficient gives, at $T_i=u^{m_i}-1$,
\[
 \sum_J A_J(T_i:i\text{ unpaired})
       \prod_{j\in J}
       \frac{C_{a_j,a_{j+1}}(T_j,T_{j+1})}{D_j}=0,
 \qquad D_j=(1+T_j)(1+T_{j+1})-\alpha.                         \tag{N.4}
\]
There are finitely many tilings. Multiply by all $D_j$, $1\le j<d$. The resulting integral analytic identity holds even at grid points where a denominator vanishes, by (N.3). Smaller source degrees vanish by choice of $d$, and larger ones have target length greater than $d$; thus no omitted source terms contribute.

The finite-weight substitution is valid exactly as in the pro-$2$ proof. At every arithmetic weight cutoff only finitely many source degrees contribute. The coefficient functions on compact orbit spaces are bounded and continuous; integration, uniform locally constant approximation, and Fubini give a continuous multiplicative map. It agrees with $\rho$ on the generators and hence on $H$.

If $\alpha=u^w$ for a positive integer $w$, restrict all test weights to exceed $w$; otherwise no exclusion is needed. This removes every possible zero denominator on the test grid. For each odd residue class $a$, the remaining tails $m=a+(p-1)p^n$ still give distinct points converging to $u^a-1$ in $p\mathbb Z_p$. Applying the analytic identity theorem successively in all variables makes the denominator-cleared version of (N.4) an identity of power series. Hence (N.4) holds in the fraction field of
\[
 R=\mathbb Z_p[[T_1,\ldots,T_d]][1/p].
\]
This ring has a uniform bound on coefficient denominators. Replacing it by $\mathbb Q_p[[T_1,\ldots,T_d]]$ would destroy the argument by allowing inverses of $D_j$ with unbounded denominators.

Choose a nonzero coefficient pattern with the maximum number of $y$ letters, and denote its tiling by $J$. Assign target types $(a_0,b_0)$ at each of its pairs and the specified source type at each singleton. This fixes an identity (N.4) with $A_J\ne0$. Consider
\[
 I_J=(D_j:j\in J)\subset R.
\]
This is prime: its disjoint equations substitute $T_{j+1}=\alpha/(1+T_j)-1$, leaving a power-series ring in the unpaired variables and one variable from each pair, with $p$ inverted. The substitutions are legitimate in the $(p,T)$ topology because their constant term is in $p\mathbb Z_p$.

For $k\notin J$, $D_k\notin I_J$. Its neighboring positions belong to different blocks; the remaining variables of these blocks are independent, and their product is not identically $\alpha$. Thus $D_k$ is a unit in $R_{I_J}$. Multiply (N.4) by $\prod_{j\in J}D_j$ and reduce modulo $I_J$ in this localization. A tiling omitting any pair of $J$ vanishes. A tiling properly containing $J$ has more $y$ letters and its coefficient is zero by maximality. Only $J$ remains:
\[
 A_J\prod_{j\in J}C_{a_0,b_0}(T_j,T_{j+1})=0
 \quad\text{in }\operatorname{Frac}(R/I_J).
\]
Every $C$ factor is a unit. The unpaired-variable ring embeds in $R/I_J$, so $A_J$ remains nonzero, a contradiction. If $J$ is empty, maximality leaves just its singleton coefficient, giving the same contradiction directly. This proves injectivity under (N.1)--(N.2).

\paragraph{The arithmetic module at $37$.}
Now set $p=37$, $F=\mathbb Q(\mu_{37})$, $K=\mathbb Q(\mu_{37^\infty})$, $P=\operatorname{Gal}(\Omega_{37}/F)$, and $H=\operatorname{Gal}(\Omega_{37}/K)$. Choose commuting lifts of $\Delta=\operatorname{Gal}(F/\mathbb Q)$ and $\Gamma=\operatorname{Gal}(K/F)$ using Sharifi's procyclic construction \cite[\S2, p.~3]{Sharifi2002}. If the chosen generator has cyclotomic value $u_0$, replace it by its $v$th power, where $v=\log(1+37)/\log(u_0)\in\mathbb Z_{37}^{\times}$. The resulting $\gamma$ still commutes with $\Delta$ and has $\chi(\gamma)=1+37$.

Iwasawa's theorem and Ferrero--Washington give freeness of $H$ \cite[Theorem~2, pp.~273--274]{Iwasawa1981}. Put $M=H^{\mathrm{ab}}$. Weak Leopoldt gives $\operatorname{rank}_{\Lambda}M=18$. Kummer duality gives that $M^+$ is $\Lambda$-torsion \cite[Theorems~3.3.3--3.3.4, Proposition~3.4.7]{SharifiIwasawaNotes}. Also $M$ is $\mathbb Z_{37}$-torsion-free because $H$ is free pro-$37$.

The minimal generators in McCallum--Sharifi, \S8, pp.~22--23, consist of $\gamma$, eighteen odd $x_m$ for $3\le m\le37$, and a single even generator $y$ of character $32$ \cite{McCallumSharifi2003}. Vandiver at $37$ and its unique irregular index $32$ were checked above. The exact sequence
\[
 0\longrightarrow M_{\Gamma}\longrightarrow P^{\mathrm{ab}}
   \longrightarrow\Gamma\longrightarrow0
\]
follows from the split semidirect product $P=H\rtimes\Gamma$. Consequently $M/(37,T)M$ has exactly eighteen odd directions and one even direction. Nakayama gives eighteen generators for $M^-$. Its rank is eighteen, since $M^+$ has rank zero. A surjection $\Lambda^{18}\to M^-$ has rank-zero kernel, which must vanish as a submodule of a free module over the domain $\Lambda$. Thus
\[
 M^-\simeq\bigoplus_{a\text{ odd}\bmod36}\Lambda e_a.           \tag{N.5}
\]

Similarly $M^+$ is cyclic, generated by the class of $y$, and is nonzero because its quotient modulo $(37,T)$ has dimension one. The explicit relation in McCallum--Sharifi, Corollary~8.3, p.~23, begins
\[
 37y-3[\gamma,y]-11[x_3,x_{29}]-[x_7,x_{25}]+\cdots=0
                                                                  \tag{N.6}
\]
in descending $p$-central degree two. Linearizing the full relator in the cyclotomic kernel and projecting onto the plus part gives an annihilator
\[
 f(T)y=0,\qquad f(T)\equiv37-3T\pmod{(37,T)^2}.                \tag{N.7}
\]
Here is the filtration justification for the coefficient congruence. Before imposing arithmetic relations, linearize the finite-rank free presentation group into $A_{\mathrm{free}}\rtimes\Gamma$, where $A_{\mathrm{free}}$ is the free $\Lambda$-module on the non-cyclotomic generators. Its filtration
\[
 (p,T)^{n-1}A_{\mathrm{free}}\rtimes p^{n-1}\Gamma
\]
respects commutators and $p$th powers; use
$(1+T)^{p^{n-1}}-1\in(p,T)^n$. Thus descending $p$-central degree three contributes only $(p,T)^2$ to each coefficient. Read the $y$ coefficient there, then pass to $M^+$. The odd commutators vanish and odd linear terms disappear in this projection.

Weierstrass preparation makes $f$ a unit times $T-\beta$, where $\beta\in37\mathbb Z_{37}$ and $v_{37}(\beta)=1$. Hence $M^+$ is a quotient of $\Lambda/(f)\simeq\mathbb Z_{37}$. It is nonzero and $\mathbb Z_{37}$-torsion-free, so
\[
 M^+\simeq\mathbb Z_{37},\qquad\gamma e_y=(1+\beta)e_y.        \tag{N.8}
\]
Equations (N.5)--(N.8) prove (N.1), with $\alpha=1+\beta$.

The lifts can preserve the preceding Frattini classes: project their abelian classes with the integral $\Delta$-idempotents and lift back through $H\to M$. The projections fix their prescribed residual characters. Each odd projected class generates its free rank-one $\Lambda$-eigenspace, and the even class generates $M^+$. A change of finite free presentation generators with the same Frattini classes is induced by an automorphism: the endomorphism is surjective by the Frattini criterion and injective by Hopficity of a finitely generated profinite group. Such a change fixes descending $p$-central degree two: both $p$th powers and commutators retain their classes. Unit rescalings within the one-dimensional eigenspaces only rescale the displayed nonzero coefficients by units. Thus these changes preserve the information used from (N.6).

\paragraph{Recovering the unit quadratic numerator.}
For the compact basis supplied by (N.1), write $\varphi$ for the $\Gamma$-automorphism of $H$ and $E=\varphi(y)y^{-\alpha}$. The group $P$ has the presentation on $\gamma,x_a,y$ with one defining relation
\[
 R=\gamma y\gamma^{-1}\varphi(y)^{-1}.
\]
To verify the presentation, assign images to these generators in any pro-$p$ target satisfying $R$. Send the compact orbit generator $x_{a,t}$ to $\gamma^t x_a\gamma^{-t}$. The compact free universal property gives a map from $H$; its $\Gamma$-compatibility is automatic on the orbits and is exactly $R$ on $y$. Hence it extends to $H\rtimes\Gamma$.

In descending $p$-central degree two the relation has class
\[
 [\gamma,y]+(1-\alpha)y-E_2,
\]
where $E_2$ is the quadratic Magnus term after folding the orbit parameters to their augmentations. Its ordered $x_a x_b$ coefficient is, up to the commutator sign convention, $C_{a,b}(0,0)$. Because there is one defining relation, any relation's degree-two initial term is a scalar multiple of this one: conjugation fixes that class and $p$th powers of the relator vanish in that degree. No assertion about freeness of the full relation module is required.

In (N.6) the coefficients of $[\gamma,y]$ and $[x_3,x_{29}]$ are both units modulo $37$. It follows that $C_{3,29}(0,0)$ is a unit. This is (N.2), with $(a_0,b_0)=(3,29)$.

\paragraph{The actual rational image and conclusion.}
Hain--Matsumoto's actual-image comparison and Brown's rational Deligne--Ihara theorem supply the dense map $H\to U(\mathbb Q_{37})$ and its completed free odd Lie algebra \cite[Proposition~8.3, Theorem~8.4]{HainMatsumoto2001}\cite[\S2.1--2.3]{Brown2019}. Arrange that the image of $\gamma$ lies in a Levi torus by successively solving in positive weights, since $u^w-1\ne0$ in $\mathbb Q_{37}$ for every $w>0$. This converges in the weight completion. The commuting finite group $\Delta$ then also lies in that torus: its possible unipotent part would centralize $\gamma$, whose unipotent centralizer is trivial by the same weight argument.

All hypotheses of the separation argument are now verified, so the actual rational action is injective on $H$. The kernel of the pro-$37$ tripod action lies in $H$ because abelianization records the cyclotomic character, and it is trivial under the induced rational action. Thus it is trivial on $P$, proving Theorem~\ref{thm:kernel37-poles}.

This proof deals with the single even arithmetic direction at $37$. It does not supply an all-irregular-primes theorem. With several even directions or higher $\Lambda$-torsion rank, the corresponding matrix and higher-order pole separation must be proved separately.
